\chapter{Metodologia}
\label{cap:metodologia}

Este capítulo descreve o desenho metodológico adotado para avaliar a modernização do módulo
de roteamento geográfico GPSR no simulador ns-3. São apresentados: a caracterização da
pesquisa; o ambiente de desenvolvimento utilizado; os critérios de seleção da implementação
do GPSR a ser modernizada; o cenário de VANET construído para os experimentos; os
protocolos comparados; as métricas de desempenho coletadas; e o planejamento estatístico dos
experimentos. O processo técnico de diagnóstico e modernização do código-fonte, por sua vez,
é detalhado no Capítulo~\ref{cap:implementacao}.

\section{Caracterização da Pesquisa}
\label{sec:caracterizacao}

Quanto à natureza, esta pesquisa é classificada como aplicada, pois tem como produto um
artefato de software — o módulo GPSR modernizado — destinado a uso prático pela comunidade
de usuários do ns-3. Quanto à abordagem, é predominantemente quantitativa, uma vez que a
avaliação do artefato produzido se dá por meio da coleta e da análise estatística de métricas
de desempenho de rede. Quanto aos procedimentos, combina uma etapa de pesquisa bibliográfica
— a Revisão Sistemática da Literatura apresentada no Capítulo~\ref{cap:rsl} — com uma etapa
de pesquisa experimental, na qual o artefato desenvolvido é avaliado em ambiente controlado
de simulação, seguindo o delineamento descrito nas seções seguintes.

\section{Ambiente de Desenvolvimento}
\label{sec:ambiente}

Os experimentos são realizados em uma versão atual e estável do ns-3, executada em ambiente
Linux, com as ferramentas de compilação necessárias à construção do simulador e de seus
módulos (compilador C++, CMake e as bibliotecas de dependência do ns-3). A coleta de
métricas é realizada por meio do módulo FlowMonitor \cite{carneiro2009flowmonitor}, nativo do
simulador, e a análise dos dados exportados é conduzida com ferramentas de processamento de
dados e geração de gráficos externas ao ns-3. O ambiente de desenvolvimento completo,
incluindo versões específicas de software e passos de instalação, é documentado em separado,
de modo a garantir a reprodutibilidade do experimento por terceiros.

\section{Seleção da Implementação do GPSR}
\label{sec:selecao-gpsr}

Como discutido na Seção~\ref{sec:roteamento-geografico}, não há, no núcleo oficial do ns-3,
um módulo de roteamento geográfico. A implementação de referência utilizada como ponto de
partida para a modernização é a documentada por \citeonline{katsaros_positionbased_ns3}, que
deu origem a diversas cópias não-oficiais mantidas pela comunidade em repositórios de código
independentes. Dentre essas cópias, a seleção da versão a ser modernizada considera os
seguintes critérios: (i) proximidade da versão do ns-3 para a qual o código foi originalmente
escrito, de modo a minimizar o número de incompatibilidades a corrigir; (ii) completude da
implementação, isto é, presença de todos os arquivos que compõem o módulo (modelo, auxiliares
de configuração e exemplos); (iii) disponibilidade de documentação mínima sobre o
funcionamento do código; e (iv) compatibilidade de licença com a licença do próprio ns-3
(GPLv2), de modo a viabilizar a disponibilização do módulo modernizado à comunidade.

\section{Cenário de VANET}
\label{sec:cenario-vanet}

O cenário experimental representa um trecho de rodovia percorrido por um conjunto de
veículos, utilizando um modelo de mobilidade nativo do ns-3, na linha do proposto por
\citeonline{arbabi2010highway}, sem dependência de simuladores de mobilidade externos,
conforme justificado na Seção~\ref{sec:mobilidade}. São definidos como parâmetros
experimentais o número de veículos, a densidade da rede (veículos por unidade de área ou de
extensão de via) e a duração de cada execução de simulação, os quais são variados de forma
controlada para permitir observar o comportamento dos protocolos avaliados sob diferentes
condições de carga e mobilidade.

\section{Protocolos Comparados}
\label{sec:protocolos-comparados}

São comparados três protocolos de roteamento, executados sobre o mesmo cenário e os mesmos
parâmetros experimentais: o módulo GPSR modernizado, objeto deste trabalho; e os protocolos
topológicos nativos do ns-3, AODV e OLSR, cuja simulação no simulador é documentada por
\citeonline{pimenta2013simulacao}. A comparação entre um protocolo geográfico e dois
protocolos topológicos, sob condições experimentais idênticas, permite isolar o efeito da
estratégia de roteamento sobre as métricas de desempenho coletadas.

\section{Métricas de Desempenho}
\label{sec:metricas}

As métricas de desempenho são coletadas automaticamente por meio do módulo FlowMonitor
\cite{carneiro2009flowmonitor} ao final de cada execução de simulação, compreendendo: a taxa
de transferência (\textit{throughput}), medida em kbps; o atraso médio fim a fim
(\textit{delay}), medido em milissegundos; a taxa de entrega de pacotes (\textit{Packet
Delivery Ratio}, PDR), medida em percentual de pacotes entregues em relação aos enviados; e o
\textit{overhead} de roteamento, medido pela razão entre pacotes de controle e pacotes de
dados efetivamente entregues. Essas métricas são escolhidas por serem as mais frequentemente
reportadas na literatura de avaliação de protocolos de roteamento em VANETs, conforme
identificado na Revisão Sistemática da Literatura (Capítulo~\ref{cap:rsl}).

\section{Planejamento dos Experimentos}
\label{sec:planejamento}

Para cada combinação de cenário e protocolo, são realizadas 30 execuções independentes da
simulação, variando-se a semente do gerador de números aleatórios do ns-3
(\texttt{RngRun}) a cada execução, de modo a capturar a variabilidade inerente aos aspectos
estocásticos do cenário (posicionamento inicial dos veículos e variações no padrão de
tráfego de rede). A mesma sequência de 30 sementes é aplicada aos três protocolos
comparados (Seção~\ref{sec:protocolos-comparados}), garantindo que as diferenças observadas
nas métricas de desempenho sejam atribuíveis à estratégia de roteamento, e não a variações
aleatórias entre execuções.

Os resultados de cada métrica são reportados como média das 30 execuções, acompanhada do
intervalo de confiança de 95\%. A significância estatística das diferenças observadas entre
o GPSR modernizado e os protocolos nativos é avaliada por meio de testes de hipótese
apropriados à distribuição dos dados coletados, cuja aplicação e resultados são apresentados
no Capítulo~\ref{cap:resultados}.